\section{Distribusi}
\begin{defn} Misalkan $\Omega$ suatu himpunan bagian terbuka tak kosong dari~$\R^n$. Suatu fungsi $f\colon \Omega \sto \R$ disebut
 \index{terintegralkan!secara lokal}%
 \index{secara lokal!terintegralkan}%
\df{terintegralkan secara lokal} jika $\int_K \abs{f}\,d\mu < \infty$ untuk setiap himpunan bagian kompak $K$ dari~$\Omega$. (Di sini $\mu$ adalah ukuran Lebesgue
berdimensi-$n$ yang biasa.) Keluarga semua fungsi tersebut pada $\Omega$ dinotasikan
 \index{L@$\locint\Omega$!fungsi yang terintegralkan secara lokal pada~$\Omega$}%
dengan~$\locint\Omega$.
\end{defn}

\begin{exam} Fungsi konstan $\vc 1$ pada garis real terintegralkan secara lokal, meskipun fungsi ini jelas tidak terintegralkan.
\end{exam}

Bahkan fungsi yang terintegralkan secara lokal tidak harus kontinu. Pada garis real, misalnya, fungsi karakteristik dari sebarang himpunan terukur Lebesgue
terintegralkan secara lokal. Bayangkan betapa memudahkannya jika kita dapat membicarakan secara bermakna \emph{turunan} dari sebarang fungsi demikian dan,
lebih baik lagi, diperbolehkan menurunkannya sesering yang kita inginkan. --- Selamat datang di dunia distribusi. Di sini
   \begin{enumerate}
       \item setiap fungsi yang terintegralkan secara lokal dapat dipandang sebagai suatu distribusi, dan
       \item setiap distribusi dapat diturunkan tak berhingga kali, dan lebih lanjut,
       \item semua aturan diferensiasi yang lazim dalam kalkulus dasar berlaku.
   \end{enumerate}

\begin{defn} Misalkan $\Omega$ suatu himpunan bagian terbuka tak kosong dari~$\R^n$. Suatu
 \index{distribusi}%
\df{distribusi} adalah fungsional linear kontinu pada ruang fungsi uji $\fml D(\Omega) = \fml C^\infty_c(\Omega)$ pada~$\Omega$.
Meskipun secara teknis fungsional demikian didefinisikan pada suatu ruang fungsi pada $\Omega$, kita akan mengikuti kebiasaan longgar yang lazim
dan menyebutnya sebagai \emph{distribusi pada~$\Omega$}. Himpunan semua distribusi pada $\Omega$, yaitu ruang dual dari $\fml D(\Omega)$,
 \index{D@$\fml D^*(\Omega)$ (ruang distribusi pada~$\Omega$)}%
akan dinotasikan dengan~$\fml D^*(\Omega)$.
\end{defn}

\begin{prop} Misalkan $\Omega$ suatu himpunan bagian terbuka dari~$\R^n$. Suatu fungsional linear $L$ pada ruang fungsi uji $\fml D(\Omega)$ pada $\Omega$ adalah
distribusi jika dan hanya jika $L(\phi_k) \sto 0$ setiap kali $(\phi_k)$ merupakan barisan fungsi uji sedemikian sehingga $\phi_k^{(\bs\alpha)} \sto 0$
secara seragam untuk setiap multi-indeks $\bs\alpha$ dan terdapat suatu himpunan kompak $K \subseteq \Omega$ yang memuat tumpuan setiap~$\phi_k$.
\end{prop}

\begin{proof} Lihat~\cite{Conway:1990}, proposisi IV.5.21 atau~\cite{Treves:1967}, proposisi 21.1(b).  \ns
\end{proof}

\begin{exam} Jika $\Omega$ suatu himpunan bagian terbuka tak kosong dari $\R^n$ dan $f$ suatu fungsi yang terintegralkan secara lokal pada $\Omega$, maka pemetaan
   \[ L_f\colon \fml D(\Omega) \sto \K \colon \phi \mapsto \int f\phi\,d\lambda \]
merupakan suatu distribusi. Distribusi demikian disebut
 \index{regular!distribusi}%
 \index{distribusi!regular}%
 \index{L@$L_f$ (distribusi regular yang bersesuaian dengan~$f$)}%
distribusi \df{regular}. Distribusi yang tidak regular disebut
 \index{singular!distribusi}%
 \index{distribusi!singular}%
distribusi \df{singular}.
\end{exam}

\emph{Catatan.} Misalkan $u$ suatu distribusi dan $\phi$ suatu fungsi uji pada $\Omega$. Notasi baku untuk $u(\phi)$ adalah
$\langle u,\phi \rangle$ dan $\langle \phi,u \rangle$. Kita juga menggunakan
 \index{<unoptop@$\tilde f$ (distribusi regular yang bersesuaian dengan~$f$)}%
notasi $\tilde f$ untuk $L_f$. Jadi, dalam contoh sebelumnya,

    \[ \langle \tilde f,\phi \rangle = \langle \phi,\tilde f \rangle
            = \tilde f(\phi) = \langle L_f,\phi \rangle
            = \langle \phi,L_f \rangle = L_f(\phi) \]
untuk $f \in \locint\Omega$.

\begin{exam} Jika $\mu$ suatu ukuran Borel regular pada himpunan bagian terbuka tak kosong $\Omega$ dari~$\R^n$ dan $f$ suatu fungsi yang terintegralkan secara lokal
terhadap ukuran tersebut pada $\Omega$, maka pemetaan
    \[ L_\mu\colon \fml D(\Omega) \sto \R\colon \phi \mapsto \int_\Omega f\phi \,d\mu \]
merupakan suatu distribusi pada $\Omega$.
\end{exam}

\begin{notn} Misalkan $a \in \R$. Definisikan $\mu_a$, yaitu
 \index{Dirac!ukuran}%
 \index{ukuran!Dirac}%
 \index{mua@$\mu_a$ (ukuran Dirac yang terkonsentrasi di~$a$)}%
\df{ukuran Dirac} yang terkonsentrasi di $a$, pada himpunan-himpunan Borel dari $\R$, dengan
   \[ \mu_a(B) = \begin{cases}  1, &\text{jika $a \in B$} \\
                                0, &\text{jika $a \notin B$}.
                 \end{cases} \]
Distribusi yang bersesuaian kita notasikan dengan $\delta_a$ dan kita sebut
 \index{Dirac!distribusi}%
 \index{distribusi!Dirac}%
 \index{delta@$\delta_a$, $\delta$ (distribusi Dirac)}%
\df{distribusi delta Dirac di $a$}. Jadi, $\delta_a = L_{\mu_a}$. Jika $a = 0$, kita menulis $\delta$ untuk $\delta_0$.
Secara historis, distribusi $\delta$ telah menyandang nama yang terkenal menyesatkan (meskipun secara teknis
tidak salah): \emph{fungsi delta Dirac}.
\end{notn}

\begin{exam} Jika $\delta_a$ adalah \emph{distribusi delta Dirac} di suatu titik $a \in \R$ dan $\phi$ suatu fungsi uji pada~$\R$,
maka $\delta_a(\phi) = \phi(a)$.
\end{exam}

\begin{notn} Misalkan $H$ menyatakan fungsi karakteristik interval $[0,\infty)$. Fungsi ini adalah
 \index{Heaviside!fungsi}%
 \index{fungsi!Heaviside}%
\df{fungsi Heaviside}. Distribusi yang bersesuaian, $\widetilde H = L_H$, adalah
 \index{Heaviside!distribusi}%
 \index{distribusi!Heaviside}%
 \index{H@$\wt H$ (distribusi Heaviside)}%
\df{distribusi Heaviside}.
\end{notn}

\begin{exam} Jika $\wt H$ adalah \emph{distribusi Heaviside} dan $\phi \in \fml D(\R)$, maka $\wt H(\phi) = \int_0^\infty \phi$.
\end{exam}

\begin{exam} Misalkan $f$ suatu fungsi yang dapat didiferensialkan pada $\R$ dan turunannya terintegralkan secara lokal. Maka
    \begin{equation}\label{deriv_reg_distr}
        L_{f\,'}(\phi) = -L_f(\phi\,')
    \end{equation}
untuk setiap $\phi \in \fml D(\R)$.
\end{exam}

\begin{proof}[\emph{Petunjuk untuk bukti}] Integrasi parsial. \ns
\end{proof}

Sekarang timbul sebuah pertanyaan penting: Bagaimana kita dapat mendefinisikan ``turunan'' suatu distribusi? Tentunya wajar jika kita menginginkan suatu
distribusi regular berperilaku hampir sama dengan fungsi yang melahirkannya. Artinya, kita ingin turunan $(L_f)'$ dari
suatu distribusi regular $L_f$ bersesuaian dengan distribusi regular yang timbul dari turunan fungsi~$f$. Dengan kata lain,
yang kita inginkan adalah $(L_f)' = L_{f\,'}$. Menurut contoh sebelumnya, ini menyarankan bahwa untuk fungsi $f$ yang dapat didiferensialkan
yang turunannya terintegralkan secara lokal, kita seharusnya mendefinisikan
   \begin{equation}\label{deriv_distr1}
        (L_f)'(\phi) = -L_f(\phi\,')
   \end{equation}
untuk setiap $\phi \in \fml D(\R)$. Pengamatan bahwa ruas kanan persamaan~\eqref{deriv_distr1} bermakna untuk
\emph{distribusi apa pun} memotivasi definisi berikut.

\begin{defn}\label{def_deriv_distr} Misalkan $u$ suatu distribusi pada himpunan bagian terbuka tak kosong $\Omega$ dari~$\R$. Kita mendefinisikan $u\,'$, yaitu
 \index{turunan!distribusi}%
 \index{distribusi!turunan}%
\df{turunan} dari $u$,
 \index{D@$Du$ (turunan suatu distribusi~$u$)}%
 \index{<unopur@$u'$ (turunan suatu distribusi~$u$)}%
dengan
    \[ u\,'(\phi) := -u(\phi\,') \]
untuk setiap fungsi uji $\phi$ pada~$\Omega$. Kita juga menggunakan notasi $Du$ untuk turunan dari~$u$.
\end{defn}

\begin{exam} \emph{Distribusi Dirac} $\delta$ adalah turunan dari \emph{distribusi Heaviside}~$\widetilde H$.
\end{exam}

\begin{exer} Misalkan $S$ suatu fungsi signum (atau tanda) pada $\R$; artinya, misalkan $S = 2H - 1$, dengan $H$ fungsi Heaviside.
Misalkan pula $g$ suatu fungsi pada $\R$ sedemikian sehingga $g\colon x \mapsto -x - 3$ untuk $x < 0$ dan $g\colon x \mapsto x + 5$ untuk $x \ge 0$.
   \begin{enumerate}[label=\rm{(\alph*)}]
     \item Carilah konstanta $\alpha$ dan $\beta$ sedemikian sehingga ${\tilde g}' = \alpha \widetilde S + \beta \delta$.
     \item Carilah suatu fungsi $a$ pada $\R$ sedemikian sehingga ${\tilde a}' = \wt S$.
   \end{enumerate}
\end{exer}

\begin{exer} Misalkan $f(x) = \left\{
                           \begin{array}{ll}
                             1 - \abs x, & \hbox{jika $\abs x \le 1$;} \\
                                 0,      & \hbox{selain itu.}
                           \end{array}
                         \right.$ Dengan menghitung kedua ruas persamaan~\eqref{deriv_reg_distr} secara terpisah, tunjukkan bahwa persamaan itu
benar secara klasik untuk setiap fungsi uji yang tumpuannya termuat dalam interval~$[-1,1]$.
\end{exer}

\begin{exer} Misalkan $f(x) = \sbsb\chi{(0,1)}$ untuk $x \in~\R$ dan $\phi \in \fml D\bigl((-1,1)\bigr)$. Bandingkan nilai $L_{f\,'}(\phi)$ dan
$\bigl(L_f\bigr)'(\phi)$.
\end{exer}

\begin{exer} Misalkan $f(x) = \left\{
                           \begin{array}{ll}
                             x, & \hbox{jika $0 \le x \le 1$;} \\
                             x + 2, & \hbox{jika $1 < x \le 2$;} \\
                             0, & \hbox{selain itu.}
                           \end{array}
                         \right.$ Bandingkan nilai $L_{f\,'}(\phi)$ dan $\bigl(L_f\bigr)'(\phi)$ untuk suatu fungsi uji~$\phi$.
\end{exer}

\begin{exer} Misalkan $h(x) = \left\{
                           \begin{array}{ll}
                             2x, & \hbox{jika $x < -1$;} \\
                             1, & \hbox{jika $-1 \le x < 1$;} \\
                             0, & \hbox{jika $x \ge 1$.}
                           \end{array}
                         \right.$ Carilah suatu fungsi $f$ yang terintegralkan secara lokal sedemikian sehingga ${\tilde f}\,' = \tilde h + 3\delta_1 - \delta_{-1}$.
\end{exer}

Definisi~\ref{def_deriv_distr} dapat diperluas ke turunan parsial dari distribusi tanpa kesulitan.

\begin{defn} Misalkan $u$ suatu distribusi pada suatu himpunan bagian terbuka tak kosong $\Omega$ dari $\R^n$ untuk suatu $n \ge 2$ dan $\bs\alpha$
suatu multi-indeks. Kita mendefinisikan
 \index{operator diferensial!yang bekerja pada distribusi}%
 \index{D@$D^{\bs\alpha}$ (notasi multi-indeks untuk diferensiasi)}%
\df{operator diferensial} $D^{\bs\alpha} = \bigl(\frac{\partial}{\partial x_1}\bigr)^{\alpha_1} \dots
\bigl(\frac{\partial}{\partial x_n}\bigr)^{\alpha_n}$ berorde $\abs{\bs\alpha}$ yang bekerja pada $u$ dengan
   \[ D^{\bs\alpha}u(\phi) := (-1)^{\abs{\bs\alpha}}u\bigl(\phi^{(\bs\alpha)}) \]
untuk setiap fungsi uji $\phi$ pada~$\Omega$.
 \index{<unopura@$u^{\bs\alpha}$ (notasi multi-indeks untuk diferensiasi)}%
Notasi alternatif untuk $D^{\bs\alpha}u$ adalah $u^{(\bs\alpha)}$.
\end{defn}

\begin{exer} Suatu \df{dipol} (dengan momen listrik 1 di titik asal pada $\R$) dapat dipandang sebagai ``limit'' ketika $\epsilon \sto 0^+$
 \index{dipol}%
dari suatu sistem $T_\epsilon$ yang terdiri atas dua muatan $-\frac1\epsilon$ dan $\frac1\epsilon$ yang masing-masing ditempatkan di $0$ dan $\epsilon$. Tunjukkan
bagaimana hal ini dapat mengarahkan kita untuk mendefinisikan dipol secara matematis sebagai $-\delta\,'$. \emph{Petunjuk.} Pandanglah $T_\epsilon$ sebagai
distribusi $\frac1\epsilon \delta_\epsilon - \frac1\epsilon \delta$.
\end{exer}
 
